%% ============================================================================
%% §4.1 Architecture Overview  --  joint
%% ============================================================================
\subsection{Architecture Overview}
\label{sec:hw:overview}

The accelerator is a hybrid CNN NPU specialised for MobileFaceNet face verification. It supports dynamically-sized $1\times1$ convolution, $3\times3$ convolution, $3\times3$ and $7\times7$ depthwise convolution, residual addition, and PReLU activation. Inputs and outputs use INT8 quantisation throughout. The same RTL is brought up along two implementation paths sharing identical functional behaviour: a Zynq UltraScale+ FPGA bitstream used for on-board verification and rapid iteration, and an Innovus place-and-route in FreePDK45 used to obtain ASIC-level area, power, and timing data. Table~\ref{tab:asic_vs_fpga} summarises the two paths against the shared design.

\begin{table}[H]
\centering
\caption{Implementation Comparison: FPGA and ASIC Tracks}
\label{tab:asic_vs_fpga}
\begin{tabularx}{\textwidth}{l X X}
\toprule
 & \textbf{FPGA track} & \textbf{ASIC track} \\
\midrule
Target              & Zynq UltraScale+ MPSoC          & FreePDK45 standard cells \\
Compute fabric      & 14$\times$16 SA + 3$\times$3 array + dw 7$\times$7 (shared RTL) & same RTL, retargeted \\
Activation memory   & BRAM (Vivado-inferred)          & OpenRAM SRAM macro \\
Weight memory       & BRAM (1\,MB)                    & combinational ROM (no usable SRAM compiler entry) \\
Param / config RAM  & BRAM                            & BRAM-style synthesised RAM \\
Achieved clock      & 150\,MHz                        & 200\,MHz \\
Cycles / inference  & 19{,}175{,}557                  & 19{,}175{,}557 \\
Latency / image     & 127.8\,ms                       & 95.88\,ms \\
Speedup vs i7-9700 (917.37\,ms) & $\times$7.18         & $\times$9.57 \\
Implementation tool & Vivado                          & Cadence Genus + Innovus \\
Verification        & on-board PS readback            & RTL + Genus syn + Innovus post-route + SDF \\
\bottomrule
\end{tabularx}
\end{table}

The two tracks are complementary. The FPGA bitstream provides hardware-in-the-loop functional validation at the system level, with the Zynq Processing System (PS) driving the NPU and reading back the 128-d embedding through AXI. The ASIC flow quantifies physical metrics (area, power, timing closure) on a standard-cell technology, with full-flow signoff via Genus synthesis, Innovus place-and-route, and SDF-annotated gate-level simulation.
